\documentclass[11pt]{article}

\usepackage[a4paper,margin=1in]{geometry}
\usepackage{amsmath,amssymb}
\usepackage{xcolor}
\usepackage{listings}
\usepackage[T1]{fontenc}
\usepackage{microtype}
\usepackage{hyperref}
\setlength{\emergencystretch}{3em}

% --- Listings setup for Lean ---
\lstdefinelanguage{Lean}{
  morekeywords={
    import, namespace, end, inductive, def, theorem, lemma, abbrev, by, intro,
    apply, exact, simp, simp_all, simp_wf, constructor, cases, rcases, with, match,
    open, have, let, refine, refine', induction, induction', noncomputable, attribute,
    instance, obtain, by_cases, by_contra, fun, variable, termination_by, decreasing_by,
    return, if, then, else, unfold, aesop, grind, convert, all_goals, generalizing, using, set
  },
  sensitive=true,
  morecomment=[l]{--},
  morecomment=[s]{/-}{-/},
  morestring=[b]"
}
\lstset{
  language=Lean,
  basicstyle=\ttfamily\small,
  keywordstyle=\color{blue!60!black},
  commentstyle=\color{green!45!black},
  stringstyle=\color{red!50!black},
  showstringspaces=false,
  breaklines=true,
  frame=single,
  columns=fullflexible,
  mathescape=true,
  % pull chunks straight from the .lean source; hide the anchor markers
  rangebeginprefix=--\ ANCHOR:\ ,
  rangeendprefix=--\ ANCHOREND:\ ,
  includerangemarker=false,
  literate=
    {Γ}{{$\Gamma$}}1 {Δ}{{$\Delta$}}1 {φ}{{$\varphi$}}1 {ψ}{{$\psi$}}1 {χ}{{$\chi$}}1
    {α}{{$\alpha$}}1 {β}{{$\beta$}}1
    {→}{{$\to$}}1 {⟶}{{$\to$}}1 {↔}{{$\leftrightarrow$}}1 {←}{{$\leftarrow$}}1
    {·}{{$\cdot$}}1 {∀}{{$\forall$}}1 {∃}{{$\exists$}}1 {∈}{{$\in$}}1 {∉}{{$\notin$}}1
    {⊆}{{$\subseteq$}}1 {∪}{{$\cup$}}1 {∧}{{$\land$}}1 {∨}{{$\lor$}}1 {¬}{{$\neg$}}1
    {⊥}{{$\bot$}}1 {⊢}{{$\vdash$}}1 {∅}{{$\emptyset$}}1 {≤}{{$\leq$}}1 {≥}{{$\geq$}}1
    {≠}{{$\neq$}}1 {~}{{$\sim$}}1 {⟨}{{$\langle$}}1 {⟩}{{$\rangle$}}1
    {‹}{{\guilsinglleft}}1 {›}{{\guilsinglright}}1
    {₁}{{$_1$}}1 {₂}{{$_2$}}1
    {á}{{\'a}}1 {…}{{$\dots$}}1 {—}{{---}}1
}

\newcommand{\Imp}{\rightarrow}
\newcommand{\Not}{\neg}
\newcommand{\Prov}{\vdash}
\newcommand{\Perm}{\sim}

\title{\textbf{AI-Reconstructed Proofs: A Reference}\\[2pt]
\large Complete Lean~4 proofs produced cold by Harmonic Aristotle, in chunks}
\author{\textbf{Rudik Rompich} \\ Universidad del Valle de Guatemala \\ \texttt{rom19857@uvg.edu.gt}}
\date{June 2026}

\begin{document}
\maketitle

\begin{abstract}
For reference, this document presents the complete Lean~4 proofs that the automated theorem-proving
agent \emph{Harmonic Aristotle} produced when asked to reconstruct the project's main results
\emph{cold} --- given only the definitions and statements, with the human proofs withheld. Rather than
a raw dump, each proof is broken into chunks, and every chunk is paired with the mathematics it
realizes. The code is pulled verbatim from the verified source files in \texttt{aristotle/}; only the
listing anchors are suppressed. Each proof was checked under Lean~v4.29.0, compiles with no
\texttt{sorry}, and depends only on \texttt{propext}, \texttt{Classical.choice}, \texttt{Quot.sound}.
How these compare to the hand proofs is discussed in the two main reports; here we record the
machine-checked evidence.
\end{abstract}

\tableofcontents
\bigskip

%==================================================================
\section{A derived rule: classical proof by cases}
Given the natural-deduction calculus and the structural lemma \texttt{weakening}, Aristotle was asked
to prove the classical \emph{proof-by-cases} rule:
\[
  \frac{\Gamma,\varphi \Prov \chi \qquad \Gamma,\Not\varphi \Prov \chi}{\Gamma \Prov \chi}.
\]
Its strategy is a classical \emph{reductio}: to derive $\chi$, assume $\Not\chi$ and reach $\bot$.
From the first branch it builds $\Gamma,\Not\chi \Prov \Not\varphi$ (if $\varphi$ held, branch~1 would
give $\chi$, contradicting $\Not\chi$); from the second, symmetrically, $\Gamma,\Not\chi \Prov
\Not\Not\varphi$. Then $\Not$-elimination on $\Not\Not\varphi$ and $\Not\varphi$ yields $\bot$.
\lstinputlisting[linerange=bycases-bycases]{../../aristotle/NDCore.lean}

%==================================================================
\section{Completeness of natural deduction}
Given the syntax, semantics, the calculus, and the elementary metatheory --- but \emph{not} Kalmár's
lemma or the discharge construction --- Aristotle reconstructed the whole completeness theorem,
rediscovering Kalmár's method under its own names.

\subsection{Signed literals and the literal context}
For a fixed valuation $v$ it introduces the \emph{signed literal} of an atom ($p$ if $v(p)$, else
$\Not p$), the finite set of atoms occurring in a formula, and the \emph{literal context}
$\Gamma_v$ collecting the signed literals over a list of atoms.
\lstinputlisting[linerange=adefs-adefs]{../../aristotle/NDFull.lean}

\subsection{Kalmár's lemma}
The key lemma: for atoms covering $\varphi$, the literal context $\Gamma_v$ \emph{decides} $\varphi$
syntactically --- it derives $\varphi$ when $v\models\varphi$ and $\Not\varphi$ otherwise. Both halves
are carried together so the induction on $\varphi$ goes through (the negation case feeds them into
each other). The statement:
\lstinputlisting[linerange=akalmarsig-akalmarsig]{../../aristotle/NDFull.lean}
\noindent and the proof, by induction on $\varphi$ (atom / negation / implication), built from the
ND constructors and closed with Mathlib automation:
\lstinputlisting[linerange=kalmarBody-kalmarBody]{../../aristotle/NDFull.lean}

\subsection{Discharging the atoms}
The literal hypotheses are removed one atom at a time: if $\varphi$ is derivable from $\Gamma_v$ for
\emph{every} valuation $v$, it is derivable from $\emptyset$. The step perturbs $v$ at a single atom
(forcing it true, then false) and combines the two derivations with the proof-by-cases rule. The
statement:
\lstinputlisting[linerange=aelimsig-aelimsig]{../../aristotle/NDFull.lean}
\noindent and the proof:
\lstinputlisting[linerange=elimBody-elimBody]{../../aristotle/NDFull.lean}

\subsection{Assembling completeness}
For a tautology $\varphi$, Kalmár's lemma gives $\Gamma_v \Prov \varphi$ for every $v$; discharge then
yields $\emptyset \Prov \varphi$.
\lstinputlisting[linerange=amain-amain]{../../aristotle/NDFull.lean}

%==================================================================
\section{Correctness of quicksort}
Given only the definition of \texttt{quicksort} and its unfolding equation --- with the permutation
and sortedness lemmas withheld --- Aristotle reconstructed correctness.

\subsection{The output is a permutation}
Here it diverged from the hand proof: instead of rearranging the partitions structurally, it reduced
$\Perm$ to equality of element \emph{multiplicities}, $a \Perm b \iff \forall x,\ \mathrm{count}(x,a)
= \mathrm{count}(x,b)$, and discharged the resulting counting goal by automation
(\texttt{simp}/\texttt{grind}).
\lstinputlisting[linerange=qaperm-qaperm]{../../aristotle/QuicksortFull.lean}

\subsection{The output is sorted}
Sortedness follows the familiar shape: by strong induction on length, both recursive results are
sorted, every element of the small half is $\leq$ the pivot and every element of the large half is
$\geq$ it, and \texttt{pairwise\_append}/\texttt{pairwise\_cons} assemble the concatenation.
\lstinputlisting[linerange=qasorted-qasorted]{../../aristotle/QuicksortFull.lean}

\subsection{Correctness}
The two lemmas conjoin into the specification of a sorting function.
\lstinputlisting[linerange=qamain-qamain]{../../aristotle/QuicksortFull.lean}

\end{document}
